% Bölüm 18 — taslak plan. Tam metin yazıldığında bu dosyanın yerini alır.
\chapter{Gödel'in Eksiklik Teoremleri ve Kanıtlanabilirlik Mantığı}\label{bol:godel}

\begin{bolumhedefleri}
\begin{itemize}
  \item Aritmetikleştirmeyi ve Gödel numaralamasını
  \item Köşegen önsavını ve birinci eksiklik teoremini
  \item İkinci eksiklik teoremini ve Löb teoremini
  \item GL mantığını ve Solovay tamlık teoremini
\end{itemize}
\end{bolumhedefleri}

\begin{taslakplan}
\textbf{Planlanan kesimler}
\begin{enumerate}[label=\thechapter.\arabic*, leftmargin=2.4em]
  \item Peano Aritmetiği ve temsil edilebilirlik
  \item Gödel numaralaması
  \item Köşegen önsavı
  \item Birinci ve ikinci eksiklik teoremleri
  \item Löb teoremi
  \item Kanıtlanabilirlik mantığı GL ve Solovay teoremi
  \item Lean 4 formalizasyonu: Foundation projesi
\end{enumerate}

\textbf{Planlanan Lean örnekleri}
\begin{itemize}
  \item \lc{Foundation kütüphanesinden teorem ifadeleri ve okunuşları}
\end{itemize}

\textbf{Temel kaynaklar:} \citet{godel1931,smith2013,boolos1993,solovay1976,saitou2026,foundation}.
\end{taslakplan}

